\scanpage{203}
En plus de la \uline{structure entière} sur le \deleted{\unsure{de}} \inserted{foncteur} de Hodge-DeRham, relatif à un pt complexe, il faut \uline{ajouter} la \uline{décomposition en types}, également de nature transcendante (elle est compatible avec la filtration canonique). On peut considérer que cette dernière \deleted{\unsure{associe}} \inserted{est déterminée via filtration} par l'\uline{opération involutive} $\tau'$ \inserted{$\unsure{\tau}$ $\frac12$ linéaire} dans $\xi_{x,\infty}$, provenant de l'écriture
\[
\xi_{x,\infty}=\xi_{x,\mathbf R}\otimes_{\mathbf R}\mathbf C
\qquad\text{où}\quad\xi_{x,\mathbf R}\simeq\xi_{x,\mathbf Z}\otimes_{\mathbf Z}\mathbf R.
\]

Si $k(a)$ est réel, on a une autre application $\frac12$ linéaire involutive $\tau''$, provenant de
\[
\xi_{x,\infty}\simeq\unsure{\xi_{a,\infty_{\mathbf R}}}\otimes_{\mathbf R}\mathbf C
\]
et $\tau$ sur $(\xi_{x,\mathbf Z})_{\mathbf C}$ n'est autre que $\tau'\tau''$
\diagram{\begin{tikzcd}[column sep=small,row sep=small]
\tau={} & \tau' \arrow[d,no head,bend right=20] & \tau'' \arrow[d,no head,bend left=20] & {}=\tau''\tau'\\
& \substack{\frac12\text{ linéaire}\\\text{\unsure{transcendant}}} & \substack{\frac12\text{ linéaire}\\\text{algébrique}} &
\end{tikzcd}}
\margin{à vérifier}

La situation locale \inserted{(dans le cas réel)} est \unsure{connue} par la connaissance de $\xi_{\unsure{x},\infty}^{\tau''}=(\xi_{a,\infty})_{\mathbf R}$ \inserted{(Hodge-DeRham réel)}, muni de $\tau=\tau'$ (involution transcendante, qui est $\mathbf R$-linéaire sur le dit espace), ou \inserted{mieux par celle de $\tau=\tau''$ opérant sur $\xi_{x,\infty}^{\tau'}\simeq\xi_{x,\mathbf R}$,} \uline{\unsure{évidemt} de $\tau$ sur $\xi_{a,\mathbf Z}$}, donc au choix de
\[
\boxed{\tau\text{ sur }\xi_{a,\ell}}
\]
[or c'est là \inserted{l'opération de} l'\unsure{élément} de $\pi_1(X)$, \uline{frobénius} en la \unsure{place} réelle -- !]
\margin{à vérifier}
\margin{202}
